%%%%%%%%%%%%%%%%%%%%%%%%%%%%%%%%%%%%%%%%%%%%%%%%%%%%%%%%%%%%%%%%%%%%%%%%%%%%%%%%%%%%%%%%%%%%%%%%%%%%%%
%
%   Filename    : chapter_2.tex 
%
%   Description : This file will contain your review of related works.
%                 
%%%%%%%%%%%%%%%%%%%%%%%%%%%%%%%%%%%%%%%%%%%%%%%%%%%%%%%%%%%%%%%%%%%%%%%%%%%%%%%%%%%%%%%%%%%%%%%%%%%%%%

\chapter{Related Works}
\label{sec:relatedworks}

\section{Test}

\section{AnitoPlume}
AnitoPlume is the first interactive, real-time 3D graphics-based simulation application explicitly tailored for visualizing the changing landscape of the Taal Volcano island in the Philippines and its corresponding eruption plumes.
Developed primarily to address the challenges of volcanic hazard and crisis communication, the simulator serves as an immersive, lightweight alternative to traditional 2D heatmaps and highly complex, non-interactive geoscientific numerical solvers (such as FALL3D and Tephra2).
While geoscientific solvers excel at quantitative predictive accuracy, they require heavy high-performance computing resources and significant processing time, making them unviable for applications requiring immediate visual feedback, interactive flexibility, and low-latency user immersion.
AnitoPlume fills this gap by merging computer graphics techniques with site-specific real-world terrain details to facilitate risk awareness and "what-if" situational modeling \cite{del2026interactive}.

\subsection{Graphics Architecture and System Performance}
To bypass the overhead, heavy abstractions, and unneeded general-purpose features of commercial game engines, AnitoPlume was custom-built using a lean, native C++ open-source API stack consisting of OpenGL for graphics rendering, Eigen for numerical linear algebra, and dearIMGUI for the user interface layout.
This custom application architecture optimizes real-time performance on consumer-grade hardware.
By keeping the system dependencies minimal, the simulator successfully achieves real-time interactivity, running at an average rendering rate of 30 frames per second (FPS) on a mid-range system with 6 GB of video RAM (VRAM) \cite{del2026interactive}.

\subsection{Usability Design and Interactive Features}
The simulator also introduced several usability improvements over previous volcanic visualization systems, including simplified camera controls, interactive parameter editing, real-time plume direction tracking, and contextual information about nearby communities affected by volcanic activity.
These additions improved both accessibility and user experience while preserving interactive performance \cite{del2026interactive}.

\subsection{System Validation and Evaluation Metrics}
To evaluate the effectiveness of the simulator, \citeauthor{del2026interactive} conducted quantitative and qualitative assessments.
Visual fidelity was measured using the Structural Similarity Index (SSIM) and Root Mean Squared Error (RMSE) by comparing simulated eruption images with real-world photographs and surveillance footage of Taal Volcano.
The study reported high visual similarity, indicating that the synthesized plume closely resembled actual volcanic eruptions.
In addition, usability testing using the System Usability Scale (SUS) demonstrated that participants generally found the simulator easy to use and suitable for interactive exploration.

\subsection{Architectural Limitations and Future Trajectories}
Despite these accomplishments, the primary contribution of AnitoPlume lies in the design of its visualization framework rather than the underlying graphics API.
The study focused on plume simulation, terrain modeling, visual fidelity, and usability evaluation, while the OpenGL implementation served primarily as the rendering platform supporting these objectives.
Consequently, the original work did not investigate how the rendering architecture itself might influence performance, resource utilization, maintainability, or future extensibility.

The authors also acknowledged several opportunities for future development, including the incorporation of more advanced rendering techniques, expanded physics-based simulations, compute shader implementations, and additional GPU-accelerated features.
These recommendations highlight the need for a rendering architecture capable of supporting increasingly sophisticated graphics workloads while maintaining interactive performance.

\section{Graphics API Architecture and Migration Implications}
The modernization of AnitoPlume involves migrating the rendering pipeline from a legazy OpenGL based implementation to the Falcor rendering framework.
Falcor utilizes DirectX12 as its underlying graphics API.
This migration introduces a fundamentally different graphics architecture characterized by explicit resource management, reduced driver overhead, and improved utilization of modern multi-core processors.
These architectural changes are particularly beneficial for AnitoPlume, whose real-time volcanic plume simulation continuously updates thousands of dynamic particle billboards while rendering large terrain models and high-resolution textures.

\subsection{OpenGL Implicit State vs. DirectX 12 Explicit State}
One of the defining characteristics of OpenGL is its use of an implicit global state machine. As described in Section 2.5 of the OpenGL specifications, the rendering state is stored within the current OpenGL context, where API commands modify persistent state that remains active until explicitly changed.
Draw calls consume the current context state, including the active shader program, bound textures, vertex buffers, depth testing configuration, blending settings, and other rendering parameters.
This design simplifies the use of the graphics library because rendering commands do not need to specify every state required for each execution.
However, it introduces additional work for the graphics driver since it must determine which states have changed since the previous draw, validate that the current combination of states is valid, and translate the API calls into GPU commands.
These operations contribute to increased CPU overhead particularly in applications with a large number of draw calls.

In contrast, DirectX 12 uses an explicit state management model. Rather than relying on a mutable global context, the application explicitly specifies the graphics pipeline configuration through Pipeline State Objects (PSOs), Root Signatures, Descriptor Heaps, and Command Lists.
A Pipeline State Object PSO encapsulates a complete graphics pipeline configuration including shaders, rasterizer state, depth and stencil settings, blending configuration, and input layout. 
The PSO is created and validated before rendering begins.
During rendering, the application simply binds the appropriate PSO instead of modifying numerous individual rendering states.

This explicit approach substantially reduces driver overhead because state validation and pipeline compilation occur during PSO creation rather than at every draw call.
As a result, the driver performs fewer runtime checks and can submit commands to the GPU more efficiently.
Furthermore, DirectX 12 enables multiple CPU threads to record command lists independently, improving scalability on modern multi-core processors.
Since each command list can be recorded independently without modifying a shared global rendering context, multiple CPU threads can generate rendering commands concurrently.
Synchronization is required only when threads access shared application-managed resources, resulting in significantly less synchronization overhead than the implicit state model of OpenGL.

\subsection{Exploding the CPU Command Recording Bottleneck}
TODO: Expound more on how openGL limits CPU scalability

\subsection{Explicit Memory Management \& Resource Binding}
TODO: Describe the differences between OpenGL's implicit GPU memory management and DirectX 12's explicit resource management model.

\subsection{Falcor Rendering Framework}
TODO: Explain how falcor abstracts the DX12 features through high level rendering interfaces.
Explain how it speeds up development process while keeping the performance optimizations

\begin{comment}
%
% IPR acknowledgement: the contents within this comment are from Ethel Ong's slides on RRL.
%
Guide on Writing your Related Works chapter
 
1. Identify the keywords with respect to your research
      One keyword = One document section
                Examples: 2.1 Story Generation Systems
			 2.2 Knowledge Representation

2.  Find references using these keywords

3.  For each of the references that you find,
        Check: Is it relevant to your research?
        Use their references to find more relevant works.

4. Identify a set of criteria for comparison.
       It will serve as a guide to help you focus on what to look for

5. Write a summary focusing on -
       What: A short description of the work
       How: A summary of the approach it utilized
       Findings: If applicable, provide the results
        Why: Relevance to your work

6. At the end of each section,  show a Table of Comparison of the related works 
   and your proposed project/system

\end{comment}
















